Results are reported at the commune (\textit{gemeente} / \textit{commune} / \textit{Gemeinde}) level
since 2014, and at the level of Electoral Cantons before that date.
To the present day, registered voter counts in Belgian elections are only
available at the level of Electoral Cantons, hence registered counts are set 
to zero for commune-level entries and captured in special units attached to the corresponding Electoral Canton.

Voters residing abroad are counted against the province where they are enrolled
as part of special `foreign affairs cantons' (\emph{canton affaires étrangères} /
\emph{speciaal kanton Buza} / \emph{Kanton Auswärtige Angelegenheiten}).

In EP elections, special votes for French-speaking lists are reported separately for Canton 
Rhode-Saint-Genèse (BE\_K00201), an area part of Flanders but with a large population of French-speaking voters.
